\documentclass[12pt, a4paper]{article}

% ─── Packages ────────────────────────────────────────────────────────────────
\usepackage[a4paper, margin=2.5cm]{geometry}
\usepackage{times}
\usepackage{setspace}
\usepackage{graphicx}
\usepackage{amsmath, amssymb}
\usepackage{booktabs}
\usepackage{longtable}
\usepackage{array}
\usepackage{multirow}
\usepackage{xcolor}
\usepackage{hyperref}
\usepackage{caption}
\usepackage{subcaption}
\usepackage{enumitem}
\usepackage{fancyhdr}
\usepackage{titlesec}
\usepackage{natbib}
\usepackage{microtype}
\usepackage{parskip}
\setcounter{tocdepth}{2}

% ─── Hyperref Setup ──────────────────────────────────────────────────────────
\hypersetup{
    colorlinks=true,
    linkcolor=black,
    citecolor=blue!60!black,
    urlcolor=blue!70!black,
    pdftitle={ELEC3612 Assignment 2 -- Music Genre Classification},
    pdfauthor={Group Report},
}

% ─── Header / Footer ─────────────────────────────────────────────────────────
\pagestyle{fancy}
\fancyhf{}
\fancyhead[L]{\small ELEC3612 -- Pattern Recognition and Machine Intelligence}
\fancyhead[R]{\small Music Genre Classification}
\fancyfoot[C]{\thepage}
\renewcommand{\headrulewidth}{0.4pt}

% ─── Section formatting ───────────────────────────────────────────────────────
\titleformat{\section}{\large\bfseries}{\thesection.}{0.6em}{}[\vspace{-0.3em}\rule{\linewidth}{0.4pt}]
\titleformat{\subsection}{\normalsize\bfseries}{\thesubsection}{0.6em}{}

% ─── Paragraph spacing ───────────────────────────────────────────────────────
\setlength{\parskip}{6pt}
\setlength{\parindent}{0pt}

% ─────────────────────────────────────────────────────────────────────────────
\begin{document}

% ─── Title Page ──────────────────────────────────────────────────────────────
\begin{titlepage}
    \centering
    \vspace*{1.5cm}

    {\Large \textbf{ELEC3612 -- Pattern Recognition and Machine Intelligence}}\\[0.4cm]
    {\large School of Electrical and Information Engineering}\\[0.2cm]
    {\large The University of Sydney}\\[1.5cm]

    \rule{\linewidth}{1pt}\\[0.5cm]
    {\LARGE \textbf{Assignment 2}}\\[0.3cm]
    {\Large \textbf{Deep Learning for Audio Understanding:}}\\[0.2cm]
    {\Large \textbf{Automatic Music Genre Classification}}\\[0.3cm]
    \rule{\linewidth}{1pt}\\[1.5cm]

    \begin{tabular}{ll}
        \textbf{Course Code:}   & ELEC3612 \\[4pt]
        \textbf{Submission:}    & \today \\
    \end{tabular}

    \vspace{2.5cm}

    {\large \textbf{Abstract}}\\[0.4cm]
    \begin{minipage}{0.88\textwidth}
    \small
    \setlength{\parskip}{4pt}
    This report presents an automatic music genre classification system trained on the     \texttt{lewtun/music\_genres} dataset from HuggingFace, comprising approximately
    25{,}000 labelled audio clips spanning 20 genre categories. After exploratory data
    analysis and justified genre filtering to 12 genres (17{,}632 clips), two complementary
    deep learning architectures are developed, trained, and compared: (1) a 2-D
    Convolutional Neural Network (CNN) operating on log-scaled Mel spectrogram images, and
    (2) a bidirectional Long Short-Term Memory network (BiLSTM) operating on sequences of
    Mel-Frequency Cepstral Coefficient (MFCC) frames. Both models are benchmarked against
    a classical Support Vector Machine (SVM) baseline. The CNN captures local spectral
    texture via four convolutional blocks; the BiLSTM models temporal evolution of musical
    structure across 130 MFCC frames. Results, per-class analysis, and a comparative
    discussion are provided across all three models.
    \end{minipage}

    \vfill
    {\small \textit{Submitted in partial fulfilment of the requirements of ELEC3612.}}
\end{titlepage}

% ─── Table of Contents ───────────────────────────────────────────────────────
\setlength{\parskip}{0pt}
\tableofcontents
\setlength{\parskip}{6pt}
\thispagestyle{empty}
\newpage

\setcounter{page}{1}

% ─────────────────────────────────────────────────────────────────────────────
\section{Introduction}
\label{sec:intro}
% ─────────────────────────────────────────────────────────────────────────────

Music genre classification is a foundational task in Music Information Retrieval (MIR) that aims to automatically assign a categorical genre label to an audio recording. The explosive growth of digital music streaming platforms has made automated genre tagging commercially essential, since manual annotation at the scale of tens of millions of tracks is infeasible.

From a signal processing perspective, audio classification is non-trivial: raw waveforms
are high-dimensional, temporally structured signals whose discriminative features
(timbre, rhythm, harmony, pitch) are distributed across both time and frequency.
Classical approaches based on handcrafted features such as Mel-Frequency Cepstral
Coefficients (MFCCs) and spectral centroid were the dominant paradigm for many years
\citep{tzanetakis2002musical}, but require domain expertise and fail to capture complex
hierarchical acoustic patterns.

Deep learning---in particular Convolutional Neural Networks (CNNs) and Recurrent Neural
Networks (RNNs)---has substantially improved audio classification performance by enabling
end-to-end feature learning directly from spectrograms or raw waveforms
\citep{choi2017convolutional}. The CNN naturally captures local spectral texture in a
Mel spectrogram image, while sequential architectures such as LSTMs capture the temporal
dynamics that distinguish genres \citep{choi2017crnn}.

This report covers the full project pipeline: \textbf{Stage~1} presents the selected
dataset, exploratory data analysis (EDA), and a justified genre filtering decision.
\textbf{Stage~2} provides a formal problem definition, motivation for deep learning, and
a literature review. \textbf{Stage~3} details the implemented architectures, training
procedure, and presents comparative results across all three models.

% ─────────────────────────────────────────────────────────────────────────────
\section{Stage 1: Dataset Selection and Description}
\label{sec:dataset}
% ─────────────────────────────────────────────────────────────────────────────

\subsection{Selected Dataset}

We have selected the \textbf{Music Genres} dataset hosted on HuggingFace Datasets
(\url{https://huggingface.co/datasets/lewtun/music\_genres}). The dataset contains
approximately \textbf{25{,}000 audio clips} labelled across \textbf{20 genre categories},
with a pre-defined train/test split provided by the source.

\subsection{Dataset Characteristics}

Table~\ref{tab:dataset_summary} summarises the key properties of the dataset.

\begin{table}[h!]
    \centering
    \caption{Summary of the Music Genres dataset properties.}
    \label{tab:dataset_summary}
    \begin{tabular}{>{\bfseries}p{4.5cm} p{8.5cm}}
        \toprule
        Property & Description \\
        \midrule
        Source & HuggingFace Datasets (\texttt{lewtun/music\_genres}) \\
        Total clips & $\approx$25{,}000 \\
        Number of genres (raw) & 20 \\
        Clip duration & $\approx$30 seconds each \\
        Data type & Audio waveform arrays with associated sampling rates \\
        Task type & Multi-class supervised classification \\
        Split & Pre-provided train/test split available; project uses custom stratified 70/15/15 split \\
        Licence & Open-access; intended for research and educational use \\
        Access method & Python \texttt{datasets} library (\texttt{load\_dataset}) \\
        \bottomrule
    \end{tabular}
\end{table}

\subsection{Exploratory Data Analysis}
\label{subsec:eda}

Before model development, we conducted the following EDA steps to build intuition and
inform preprocessing decisions:

\begin{enumerate}[noitemsep, topsep=4pt]
    \item \textbf{Genre distribution:} Plotting the clip count per genre to identify class
          imbalance across the 20 categories (Figure~\ref{fig:genre_dist}).
    \item \textbf{Genre filtering and re-distribution:} After removing 8 genres, the
          filtered class distribution was re-plotted and class imbalance quantified
          (Figure~\ref{fig:genre_dist_filtered}, Figure~\ref{fig:class_imbalance}).
    \item \textbf{Waveform visualisation:} Raw audio waveforms were plotted for
          Electronic, Classical, and Hip-Hop to compare amplitude envelope and dynamic
          range across genres (Figure~\ref{fig:waveforms}).
    \item \textbf{Mel spectrogram visualisation:} Log-scaled Mel spectrograms were
          plotted for the same three genres to build intuition about inter-class
          time-frequency differences (Figure~\ref{fig:spectrograms_sample}).
    \item \textbf{Mel spectrogram grid:} One example spectrogram per retained genre was
          arranged in a 4$\times$3 grid to visualise the full 12-class problem
          (Figure~\ref{fig:spectrograms_all}).
\end{enumerate}

\begin{figure}[h!]
    \centering
    \includegraphics[width=0.85\textwidth]{figures/genre_distribution.png}
    \caption{Sample count per genre in the raw dataset (20 genres, 24{,}985 clips). The
             \textit{Unknown} category alone accounts for 20.3\% of all samples.}
    \label{fig:genre_dist}
\end{figure}

\begin{figure}[b!]
    \centering
    \begin{subfigure}[b]{0.48\textwidth}
        \includegraphics[width=\textwidth]{figures/genre_distribution_filtered.png}
        \caption{Filtered distribution (12 genres).}
        \label{fig:genre_dist_filtered}
    \end{subfigure}
    \hfill
    \begin{subfigure}[b]{0.48\textwidth}
        \includegraphics[width=\textwidth]{figures/class_imbalance.png}
        \caption{Class imbalance — dashed line marks the mean count.}
        \label{fig:class_imbalance}
    \end{subfigure}
    \caption{Genre distribution and class imbalance after filtering to 12 genres
             (17{,}632 clips). Rock and Electronic are the largest classes; Country
             is the smallest.}
\end{figure}

\begin{figure}[b!]
    \centering
    \includegraphics[width=0.9\textwidth]{figures/waveforms_sample.png}
    \caption{Raw waveforms for Electronic, Classical, and Hip-Hop. Classical exhibits
             wide dynamic range; Electronic maintains high amplitude; Hip-Hop shows
             percussive transients.}
    \label{fig:waveforms}
\end{figure}

\begin{figure}[b!]
    \centering
    \includegraphics[width=0.9\textwidth]{figures/spectrograms_sample.png}
    \caption{Log-scaled Mel spectrograms for Electronic, Classical, and Hip-Hop.
             Classical has sparse harmonic content; Electronic shows dense broadband
             energy; Hip-Hop has strong low-frequency periodicity.}
    \label{fig:spectrograms_sample}
\end{figure}

\begin{figure}[h!]
    \centering
    \includegraphics[width=\textwidth]{figures/spectrograms_all_genres.png}
    \caption{One Mel spectrogram per retained genre arranged in a 4$\times$3 grid.
             Classical and Chiptune\,/\,Glitch are visually most distinct; Rock, Punk,
             and Hip-Hop share dense, uniform energy distributions.}
    \label{fig:spectrograms_all}
\end{figure}

\subsection{Genre Filtering}
\label{subsec:filtering}

Following EDA, eight of the 20 raw genre labels were excluded from the
final classification task on the following grounds:

\begin{itemize}[noitemsep, topsep=4pt]
    \item \textbf{Unknown} --- the single largest category (20.3\% of
          all clips) carries no meaningful musical label and would
          introduce systematic noise into the target variable.
    \item \textbf{Easy Listening} --- only 13 clips; statistically
          unusable for training or evaluation.
    \item \textbf{Blues} --- only 58 clips; insufficient for reliable
          generalisation.
    \item \textbf{Soul-RnB} --- only 94 clips; same reasoning as Blues.
    \item \textbf{Spoken} --- not a musical genre; acoustically distinct
          from all other classes in a way that trivialises classification
          rather than testing genre-level discrimination.
    \item \textbf{Old-Time / Historic} --- an ill-defined catch-all
          category with vague and inconsistent labelling criteria, likely
          to introduce noise into the target variable.
    \item \textbf{Ambient Electronic} --- substantial acoustic overlap
          with the broader Electronic category, leading to ambiguous
          ground-truth labels and expected high inter-class confusion.
    \item \textbf{International} --- another ill-defined catch-all with
          no consistent acoustic identity across clips.
\end{itemize}

These exclusions are modelling decisions intended to reduce label ambiguity, extreme
class imbalance, and non-musical categories; they do not imply that the removed labels
are invalid music categories.
The remaining \textbf{12 genres} form the final classification problem
used in all subsequent modelling stages:
\textit{Chiptune\,/\,Glitch}, \textit{Classical}, \textit{Country},
\textit{Electronic}, \textit{Experimental}, \textit{Folk},
\textit{Hip-Hop}, \textit{Instrumental}, \textit{Jazz}, \textit{Pop},
\textit{Punk}, and \textit{Rock}.

After filtering, the dataset comprises \textbf{17{,}632 clips}.
Class imbalance remains present: Rock and Electronic each contribute
approximately 17.5\% of samples, while Country accounts for only 0.8\%.
This imbalance is addressed during training via class-weighted
cross-entropy loss (Section~\ref{sec:methodology}).

\subsection{Justification for Dataset Selection}

The Music Genres dataset was selected on the following criteria:

\begin{enumerate}[noitemsep, topsep=2pt]
    \item \textbf{Clear task alignment:} The dataset was purpose-built for audio
          classification, mapping directly to a well-defined supervised learning problem.
    \item \textbf{Scale:} $\approx$25{,}000 clips provides sufficient data to train
          CNN and LSTM architectures without relying entirely on transfer learning.
    \item \textbf{Practical relevance:} Music genre classification has direct real-world
          applications in recommendation systems, content tagging, and streaming services.
    \item \textbf{Deep learning compatibility:} Raw waveform audio is naturally suited to
          CNN (spectrogram images) and LSTM (MFCC frame sequences) pipelines.
    \item \textbf{Open access:} Freely available via HuggingFace \texttt{datasets},
          ensuring reproducibility for all group members.
\end{enumerate}

% ─────────────────────────────────────────────────────────────────────────────
\section{Stage 2: Problem Definition}
\label{sec:problem}
% ─────────────────────────────────────────────────────────────────────────────

\subsection{Problem Statement}

Given a raw audio recording $\mathbf{x} \in \mathbb{R}^{T}$ of duration $T$ samples
(sampled at a fixed rate $f_s$), the goal is to learn a classifier $f_\theta$ that
maps the recording to one of $K$ discrete genre labels $y \in \{1, 2, \ldots, K\}$
(where $K = 12$ after filtering):

\begin{equation}
    f_\theta(\mathbf{x}) = \hat{y} \in \{1, \ldots, K\}
    \label{eq:classifier}
\end{equation}

where $\theta$ denotes the learnable parameters of the model and $\hat{y}$ is the
predicted genre label. The model is trained to minimise the categorical cross-entropy
loss with optional class weighting to account for genre imbalance:

\begin{equation}
    \mathcal{L}(\theta) = -\frac{1}{N} \sum_{i=1}^{N} \sum_{k=1}^{K}
    w_k \, y_{ik} \log \hat{p}_{ik}
    \label{eq:loss}
\end{equation}

where $N$ is the number of training samples, $w_k$ is an optional class weight
inversely proportional to class frequency, $y_{ik}$ is the one-hot encoded ground truth,
and $\hat{p}_{ik} = \text{softmax}(f_\theta(\mathbf{x}_i))_k$ is the predicted
probability for genre $k$.

\subsection{Input Representations}

Two complementary representations are derived from each raw audio clip, one for each
model.

\subsubsection{Mel Spectrogram (CNN Input)}

Each audio clip is converted to a \textbf{log-scaled Mel spectrogram} $\mathbf{S}
\in \mathbb{R}^{128 \times 128}$, which treats the audio as a 2-D time-frequency
image suitable for convolutional processing:

\begin{equation}
    \mathbf{S}[m, t] = \log\!\left(1 + \mathbf{M}_m \cdot |\text{STFT}(\mathbf{x})[t]|^2\right)
    \label{eq:melspec}
\end{equation}

where $\mathbf{M}_m$ is the $m$-th Mel filterbank vector and $\text{STFT}(\mathbf{x})[t]$
is the Short-Time Fourier Transform at time frame $t$. The Mel scale approximates
the non-linear frequency resolution of the human auditory system and has been shown to
be highly effective for audio classification \citep{hershey2017cnn}.

\subsubsection{MFCC Sequence (LSTM Input)}

Each audio clip is also represented as a sequence of \textbf{MFCC frames} $\mathbf{C}
\in \mathbb{R}^{T \times 40}$, where $T$ is the full frame sequence length before
temporal resizing and 40 cepstral coefficients are retained per frame. The sequence is
later interpolated to $T'=130$ frames for the LSTM input. This sequential representation
preserves the temporal structure of the audio signal and serves as input to the LSTM
model.

\subsection{Target Variable and Output Space}

The target variable $y$ is a categorical label indicating the music genre of each audio
clip. This defines the problem as a \textbf{supervised multi-class single-label
classification} task. All 12 retained classes are mutually exclusive.

\subsection{Evaluation Metrics}

The following evaluation metrics will be used (Table~\ref{tab:metrics}):

\begin{table}[h!]
    \centering
    \caption{Evaluation metrics for model assessment.}
    \label{tab:metrics}
    \begin{tabular}{>{\bfseries}p{4cm} p{9cm}}
        \toprule
        Metric & Rationale \\
        \midrule
        Accuracy & Proportion of correctly classified samples; primary metric. \\[4pt]
        Weighted F1-Score & Harmonic mean of precision and recall weighted by class
                            support; robust to class imbalance. \\[4pt]
        Macro F1 & Unweighted average F1 across all classes; penalises models that
                   ignore minority genres equally. \\[4pt]
        Per-class F1 & Per-genre breakdown revealing which classes are hardest to
                       classify. \\[4pt]
        Confusion Matrix & Reveals systematic misclassifications between acoustically
                           similar genres. \\[4pt]
        Training Curves & Loss and accuracy per epoch for both models; used to diagnose
                          overfitting and convergence behaviour. \\[4pt]
        ROC-AUC (macro) & Area under the macro-averaged one-vs-rest ROC curve; measures
                          class separability independently of the decision threshold. \\
        \bottomrule
    \end{tabular}
\end{table}

\subsection{Data Splitting Strategy}

Rather than using the pre-provided HuggingFace train/test split directly,
we performed a stratified re-split of the full filtered dataset to ensure
consistent genre proportions across all three partitions. The split
ratios are:

\begin{itemize}[noitemsep, topsep=4pt]
    \item \textbf{Training set} (70\%, 12{,}342 clips): used to optimise
          model parameters $\theta$ via stochastic gradient descent.
    \item \textbf{Validation set} (15\%, 2{,}645 clips): used for
          hyperparameter tuning and early stopping.
    \item \textbf{Test set} (15\%, 2{,}645 clips): held out until
          final evaluation.
\end{itemize}

Stratified splitting preserves proportional genre representation across
all three partitions. Normalisation statistics (mean and standard deviation) are
computed on the training split only and applied to validation and test
sets to prevent data leakage.

\subsection{Why Deep Learning is the Appropriate Modelling Choice}

Several properties of audio genre classification motivate deep learning over classical
machine learning:

\begin{enumerate}[topsep=4pt]
    \item \textbf{Hierarchical feature learning.}
    Music genre is characterised by hierarchical acoustic patterns---from low-level
    spectral energy to mid-level rhythmic motifs to high-level stylistic identity.
    Deep networks learn these hierarchical representations automatically without manual
    feature engineering \citep{lecun2015deep}.

    \item \textbf{Temporal modelling.}
    Musical structure unfolds over time. LSTMs are designed to model long-range temporal
    dependencies and are particularly effective for genres where rhythmic and harmonic
    development over the clip duration is the primary discriminative signal
    \citep{choi2017crnn}.

    \item \textbf{High-dimensional input space.}
    A 30-second clip at 22{,}050 Hz contains over 660{,}000 samples. Classical methods
    (SVM, Random Forest) require explicit dimensionality reduction or handcrafted feature
    engineering before they can be applied effectively. Deep networks learn compact,
    discriminative representations directly from high-dimensional inputs.

    \item \textbf{State-of-the-art performance.}
    The current state of the art on standard audio classification benchmarks is held
    exclusively by deep learning models. On the GTZAN genre dataset, deep CNN-based
    models consistently outperform classical SVM+MFCC baselines by a substantial margin
    \citep{choi2017convolutional}.
\end{enumerate}

% ─────────────────────────────────────────────────────────────────────────────
\section{Related Work}
\label{sec:litreview}
% ─────────────────────────────────────────────────────────────────────────────

\citet{tzanetakis2002musical} established the seminal framework for music genre
classification using handcrafted features---MFCCs, spectral centroid, spectral rolloff,
and beat histograms---fed into Gaussian mixture model and $k$-nearest neighbour
classifiers on the now-standard GTZAN dataset. While highly influential, performance was
limited to around 61\% accuracy on a 10-class problem, constrained by the expressiveness
of handcrafted features. Their MFCC feature extraction pipeline forms the basis of our
classical SVM baseline.

\citet{choi2017convolutional} demonstrated that CNNs operating on log-amplitude
spectrograms achieve strong performance on music tagging, and studied how pooling
strategies along the time and frequency axes affect the ability to capture rhythmic
versus spectral patterns. This work directly motivates our CNN 2-D architecture
operating on Mel spectrogram images.

\citet{choi2017crnn} proposed the Convolutional Recurrent Neural Network (CRNN),
combining a CNN for local spectral feature extraction with an RNN for temporal sequence
modelling---achieving state-of-the-art results on genre and mood tagging benchmarks.
This architecture directly motivates our LSTM model: rather than implementing a full
CRNN, we isolate the recurrent component by feeding MFCC sequences directly into a
bidirectional LSTM, allowing us to compare the contribution of temporal modelling
independently of convolutional feature extraction.

Table~\ref{tab:litreview} provides a concise summary of these three works and their
connection to our approach.

\begin{table}[h!]
    \centering
    \caption{Summary of related work and connection to our methodology.}
    \label{tab:litreview}
    \begin{tabular}{p{3.2cm} p{3.0cm} p{3.0cm} p{3.4cm}}
        \toprule
        \textbf{Reference} & \textbf{Model} & \textbf{Input} & \textbf{Connection to our work} \\
        \midrule
        \citet{tzanetakis2002musical} & GMM / $k$-NN & MFCCs, ZCR, beat histogram &
            Foundational framework; our SVM baseline replicates this feature pipeline \\[6pt]
        \citet{choi2017convolutional} & CNN & Log spectrogram &
            Directly motivates our CNN 2-D architecture on Mel spectrograms \\[6pt]
        \citet{choi2017crnn} & CRNN & Mel spectrogram &
            Motivates temporal modelling; we isolate the LSTM component on MFCC sequences \\
        \bottomrule
    \end{tabular}
\end{table}

% ─────────────────────────────────────────────────────────────────────────────
\section{Methodology}
\label{sec:methodology}
% ─────────────────────────────────────────────────────────────────────────────

\subsection{Preprocessing Pipeline}

Two separate feature extraction pipelines are applied---one per model---after a common
set of audio normalisation steps.

\subsubsection{Common Audio Normalisation}
\begin{itemize}[noitemsep, topsep=2pt]
    \item Resample all clips to a uniform sampling rate of $f_s = 22{,}050$~Hz.
    \item Convert stereo recordings to mono by averaging channels.
    \item Pad or trim clips to a fixed duration of 30 seconds.
\end{itemize}

\subsubsection{CNN Branch: Mel Spectrogram Extraction}
\begin{itemize}[noitemsep, topsep=2pt]
    \item Compute Mel spectrogram with 128 Mel filterbanks, a window size of 2{,}048
          samples, and a hop length of 512 samples.
    \item Apply log-amplitude compression: $\log(1 + S)$.
    \item Resize to a fixed $128 \times 128$ image.
    \item Normalise using training-set mean and standard deviation (zero mean, unit standard deviation).
    \item Save to disk as \texttt{spectrograms\_\{train,val,test\}.pt}.
\end{itemize}

\subsubsection{LSTM Branch: MFCC Extraction}
\begin{itemize}[noitemsep, topsep=2pt]
    \item Extract 40 MFCC coefficients per frame using a hop length of 512 samples.
          At 22{,}050~Hz, a 30-second clip produces $\lfloor 661{,}500 / 512 \rfloor + 1
          \approx 1{,}292$ frames. This sequence is then compressed to a fixed length of
          $T' = 130$ frames via temporal interpolation, ensuring a uniform input shape
          across all clips.
    \item Normalise each of the 40 coefficient dimensions to zero mean and unit standard deviation using training-set mean and standard deviation.
    \item Save to disk as \texttt{mfccs\_\{train,val,test\}.pt} with shape
          $(N, 130, 40)$.
\end{itemize}

Preprocessing is performed once offline; tensors are loaded directly during training to
avoid repeated computation.

\subsection{Model 1: 2-D Convolutional Neural Network}
\label{subsec:cnn}

\textbf{Input:} Log-scaled Mel spectrogram tensor of shape $(1, 128, 128)$.

\textbf{Architecture:} Four convolutional blocks with filter counts 32, 64, 128, and
256, each following the pattern:
\[
\texttt{Conv2D}(3{\times}3,\,\text{pad}=1) \to \texttt{BatchNorm2D} \to \texttt{ReLU} \to \texttt{MaxPool2D}(2{\times}2)
\]
After the four blocks, an \texttt{AdaptiveAvgPool2D}$(4,4)$ fixes the spatial size
regardless of residual dimension variation, followed by:
\[
\texttt{Flatten} \to \texttt{Dropout}(0.5) \to \texttt{Linear}(4096, 512) \to \texttt{ReLU} \to \texttt{Dropout}(0.3) \to \texttt{Linear}(512, 12)
\]
The final layer outputs raw logits for $K=12$ classes.

\textbf{Training:} Adam optimiser ($\text{lr}=10^{-3}$, weight decay $10^{-4}$);
\texttt{CrossEntropyLoss} with inverse-frequency class weights; \texttt{ReduceLROnPlateau}
scheduler (factor 0.5, patience 5); gradient clipping (\texttt{max\_norm}=1.0); early
stopping on validation accuracy (patience 10). Best checkpoint saved as
\texttt{best\_cnn.pth}.

\subsection{Model 2: Bidirectional LSTM}
\label{subsec:lstm}

\textbf{Input:} MFCC sequence tensor of shape $(130, 40)$.

\textbf{Architecture:} A 2-layer bidirectional LSTM with \texttt{hidden\_size}=256,
\texttt{dropout}=0.3 between layers. Processing the sequence in both temporal directions
allows the model to incorporate past and future context at each time step. The final
hidden states of the forward and backward passes at the last layer ---
$\mathbf{h}_{-2}$ and $\mathbf{h}_{-1}$ --- are concatenated to form a 512-dimensional
context vector, followed by:
\[
\texttt{Dropout}(0.5) \to \texttt{Linear}(512, 256) \to \texttt{ReLU} \to \texttt{Dropout}(0.3) \to \texttt{Linear}(256, 12)
\]

\textbf{Training:} Identical setup to the CNN --- Adam ($\text{lr}=10^{-3}$, weight
decay $10^{-4}$), class-weighted \texttt{CrossEntropyLoss}, \texttt{ReduceLROnPlateau},
gradient clipping (\texttt{max\_norm}=1.0), and early stopping (patience 10). Best
checkpoint saved as \texttt{best\_lstm.pth}.

\subsection{Training Configuration and Hyperparameter Selection}
\label{subsec:hyperparams}

Hyperparameters including learning rate ($10^{-3}$), batch size (64), dropout rates
(0.5 and 0.3), and weight decay ($10^{-4}$) were selected based on established
practices for deep learning on audio classification tasks---specifically, Adam's
recommended default learning rate \citep{kingma2014adam} and dropout values consistent
with those reported in \citet{srivastava2014dropout}, where a rate of 0.5 is recommended
for hidden layers. The learning rate was further adapted during training using a
\texttt{ReduceLROnPlateau} scheduler (factor 0.5, patience 5), which automatically
reduced the learning rate when validation loss plateaued, providing adaptive
optimisation without manual intervention between runs.

\subsection{Classical Baseline}

An SVM classifier with an RBF kernel is trained on per-clip MFCC summary statistics,
following the feature extraction approach of \citet{tzanetakis2002musical}. For each
clip, the mean and standard deviation of the MFCC sequence are computed across the time
axis, producing a $2 \times 40 = 80$-dimensional feature vector. Features are
standardised with \texttt{StandardScaler} (fit on the training set only), and the
classifier is trained with \texttt{SVC(kernel="rbf", C=10, gamma="scale",
class\_weight="balanced")}. This baseline quantifies the performance gain from
end-to-end deep feature learning over handcrafted representations.

% ─────────────────────────────────────────────────────────────────────────────
\section{Results}
\label{sec:results}
% ─────────────────────────────────────────────────────────────────────────────

\subsection{Overall Model Comparison}
\label{subsec:overall}

Table~\ref{tab:overall_results} summarises test-set performance across all three models.
The CNN achieves the highest accuracy (50.59\%) and weighted F1 (0.507), outperforming
the SVM baseline by 2.5 percentage points and the BiLSTM by 10.2 percentage points.
Notably, all three models exceed random-chance performance (8.3\% for 12 classes) by a
substantial margin, and the deep learning models both outperform the classical baseline,
confirming the value of end-to-end feature learning for this task. The relatively modest
absolute accuracy across all models reflects the inherent difficulty of genre
classification on a 12-class problem with acoustically overlapping categories.

\begin{table}[h!]
    \centering
    \caption{Test-set performance comparison across all three models (2{,}645 samples).}
    \label{tab:overall_results}
    \begin{tabular}{lcccc}
        \toprule
        \textbf{Model} & \textbf{Accuracy} & \textbf{Weighted F1} & \textbf{Macro F1} & \textbf{Macro AUC} \\
        \midrule
        CNN (Mel spectrogram)   & \textbf{50.59\%} & \textbf{0.507} & \textbf{0.487} & \textbf{0.893} \\
        SVM baseline (MFCC)     & 48.05\%          & 0.475          & 0.437          & 0.865          \\
        BiLSTM (MFCC sequence)  & 40.42\%          & 0.400          & 0.351          & 0.789          \\
        \bottomrule
    \end{tabular}
\end{table}

\subsection{Per-Class F1 Scores}

Table~\ref{tab:perclass_f1} provides the per-genre F1 score for each model. Hip-Hop
and Classical are consistently the strongest performers across all models; Jazz and Pop
are consistently the weakest.

\begin{table}[h!]
    \centering
    \caption{Per-class F1 scores on the test set for all three models.}
    \label{tab:perclass_f1}
    \begin{tabular}{lccc}
        \toprule
        \textbf{Genre} & \textbf{CNN} & \textbf{SVM} & \textbf{BiLSTM} \\
        \midrule
        Chiptune / Glitch  & 0.560 & 0.449 & 0.365 \\
        Classical          & 0.714 & 0.704 & 0.594 \\
        Country            & 0.444 & 0.368 & 0.111 \\
        Electronic         & 0.486 & 0.499 & 0.429 \\
        Experimental       & 0.478 & 0.425 & 0.361 \\
        Folk               & 0.596 & 0.494 & 0.440 \\
        Hip-Hop            & \textbf{0.762} & 0.609 & 0.526 \\
        Instrumental       & 0.354 & 0.372 & 0.274 \\
        Jazz               & 0.265 & 0.079 & 0.088 \\
        Pop                & 0.168 & 0.213 & 0.156 \\
        Punk               & 0.601 & \textbf{0.609} & 0.516 \\
        Rock               & 0.412 & 0.420 & 0.352 \\
        \bottomrule
    \end{tabular}
\end{table}

\subsection{Confusion Matrices}

Figures~\ref{fig:cm_cnn}--\ref{fig:cm_svm} show row-normalised confusion matrices for
each model, revealing systematic misclassification patterns.

\begin{figure}[htbp]
    \centering
    \begin{subfigure}[b]{0.48\textwidth}
        \includegraphics[width=\textwidth]{figures/cnn_confusion_matrix.png}
        \caption{CNN}
        \label{fig:cm_cnn}
    \end{subfigure}
    \hfill
    \begin{subfigure}[b]{0.48\textwidth}
        \includegraphics[width=\textwidth]{figures/lstm_confusion_matrix.png}
        \caption{BiLSTM}
        \label{fig:cm_lstm}
    \end{subfigure}

    \vspace{0.8em}
    \begin{subfigure}[b]{0.56\textwidth}
        \centering
        \includegraphics[width=\textwidth]{figures/svm_confusion_matrix.png}
        \caption{SVM baseline}
        \label{fig:cm_svm}
    \end{subfigure}
    \caption{Row-normalised confusion matrices for all three models. Strong diagonals
             for Hip-Hop and Classical across all models; Pop, Jazz, and Country show
             the most off-diagonal mass.}
    \label{fig:cm_all}
\end{figure}

\subsubsection*{CNN Confusion Analysis}

The CNN confusion matrix (Figure~\ref{fig:cm_cnn}) reveals several acoustically
motivated error patterns. The strongest off-diagonal entry is
Chiptune\,/\,Glitch$\to$Electronic (24.9\%), which is acoustically expected: both
genres share synthesised timbres, dense mid-to-high frequency energy, and the absence
of acoustic instrumentation. The reverse confusion (Electronic$\to$Chiptune, 15.0\%)
confirms this is a bidirectional boundary failure driven by shared production
characteristics rather than model bias. Country is predicted as Rock in 27.3\% of
cases, reflecting the shared acoustic and electric guitar texture that defines both
genres; without lyrical content as a signal, the spectrogram representations are
near-indistinguishable for these clips. Classical is misclassified as Instrumental in
approximately 23\% of cases --- clips that lack recognisable orchestral markers
(e.g.\ solo piano passages or sparse chamber pieces) are indistinguishable from the
broadly-defined Instrumental class. Jazz achieves only 23.9\% diagonal accuracy,
scattering primarily into Experimental ($\sim$17\%) and Instrumental ($\sim$15\%):
Jazz's complex, non-periodic harmonic structure and low class support (46 test samples)
combine to make it the hardest class for the CNN. Pop is correctly identified only
$\sim$15\% of the time, with $\sim$20\% predicted as Rock --- a direct consequence of
Pop's stylistic heterogeneity and heavy borrowing of rock instrumentation.

\subsubsection*{BiLSTM Confusion Analysis}

The BiLSTM matrix (Figure~\ref{fig:cm_lstm}) shows substantially more diffuse
off-diagonal mass, consistent with the model's lower overall accuracy. The most striking
failure is Country, where only $\sim$9\% of clips are correctly classified and $\sim$32\%
are predicted as Rock. This is worse than the CNN's already-poor Country performance,
suggesting that the temporal evolution of MFCC coefficients provides almost no signal to
distinguish Country from Rock at the frame sequence level. Notably, the BiLSTM
outperforms the CNN on Classical (66.2\% vs.\ 60.8\%), which is the one genre where
temporal structure is genuinely informative: Classical music has distinctive long-range
harmonic development and dynamic variation that unfolds over time, rewarding the
recurrent model's sequential processing. Jazz is near-unclassifiable for the BiLSTM
($<$9\% diagonal), with confusion spread broadly across Instrumental, Experimental,
and Electronic --- the MFCC sequence representation provides insufficient structure to
capture Jazz's defining characteristics of improvisation and complex chord voicings.

\subsubsection*{SVM Confusion Analysis}

The SVM matrix (Figure~\ref{fig:cm_svm}) is notable for its strong Classical
performance: 77.0\% correct, the highest of any model on any single genre. This
reflects the fact that MFCC statistics (mean and standard deviation across the clip)
effectively encode Classical music's sparse, structured harmonic content and wide
dynamic range --- properties that aggregate well into a fixed-length feature vector.
Hip-Hop (61.4\%) is similarly well-handled, as its characteristic bass emphasis and
rhythmic periodicity are faithfully captured in MFCC summary statistics. The SVM
degrades most sharply on Jazz, where the broad confusion pattern mirrors both deep
models, confirming that Jazz's difficulty is a property of the acoustic content and
class size rather than any particular model architecture. Pop and Country remain
difficult for the same reasons as in the deep models, with Pop scattering across
Rock, Punk, and Electronic, and Country confusing primarily with Rock.

\subsection{Training Curves}

Figures~\ref{fig:curves_cnn} and~\ref{fig:curves_lstm} show the training and validation
loss and accuracy per epoch for the CNN and BiLSTM respectively.

\begin{figure}[h!]
    \centering
    \includegraphics[width=0.9\textwidth]{figures/cnn_training_curves.png}
    \caption{CNN training and validation loss (left) and accuracy (right) per epoch.
             Best checkpoint saved at epoch 35. Training accuracy continues rising while
             validation accuracy plateaus around epoch 20, indicating moderate overfitting.}
    \label{fig:curves_cnn}
\end{figure}

\begin{figure}[h!]
    \centering
    \includegraphics[width=0.9\textwidth]{figures/lstm_training_curves.png}
    \caption{BiLSTM training and validation loss (left) and accuracy (right) per epoch.
             Best checkpoint saved at epoch 21. Validation loss diverges sharply after
             epoch 5 while training loss continues to fall, a textbook overfitting signature.}
    \label{fig:curves_lstm}
\end{figure}

The CNN training curves (Figure~\ref{fig:curves_cnn}) show well-behaved convergence: both
training and validation loss decrease consistently in early epochs, with the best
checkpoint saved at epoch~35. From around epoch~20 onward, training accuracy continues
to climb (reaching $\sim$63\%) while validation accuracy plateaus at $\sim$55\%,
indicating moderate overfitting. The validation accuracy curve is noticeably noisy
throughout, reflecting the inherent difficulty and class overlap of the 12-genre
problem rather than training instability. The train-val gap is manageable and was
partially controlled by the two dropout layers (0.5 and 0.3) and early stopping.

The BiLSTM training curves (Figure~\ref{fig:curves_lstm}) tell a markedly different
story. Validation loss decreases only until approximately epoch~5, after which it
diverges sharply --- rising from $\sim$1.8 to nearly 4.0 by epoch~21 --- while training
loss continues falling to $\sim$0.25. This divergence is a textbook overfitting
signature: the model memorises the training set rather than learning generalisable
temporal patterns from the MFCC sequences. The resulting train-val accuracy gap
($\sim$82\% vs.\ $\sim$43\% at epoch~21) is far larger than the CNN's, confirming that
the recurrent architecture overfits much more aggressively on this dataset. Early
stopping correctly saved the best checkpoint at epoch~21 before the overfitting became
catastrophic.

\subsection{ROC-AUC Analysis}
\label{subsec:roc}

Figure~\ref{fig:roc} shows the macro-averaged one-vs-rest ROC curves for all three
models. Each curve is computed by averaging the per-class ROC curves across all 12
genres, with the area under the mean curve (macro AUC) reported in the legend and in
Table~\ref{tab:overall_results}. A macro AUC of 1.0 represents perfect discrimination;
0.5 corresponds to a random classifier. The CNN achieves the highest macro AUC (0.893),
consistent with its accuracy and F1 superiority. Notably, the SVM's AUC (0.865) is
substantially closer to the CNN's than their 2.5 percentage point accuracy gap would
suggest --- indicating that the SVM's probability estimates are well-calibrated and that
handcrafted MFCC statistics provide strong class-separability, even if the decision
boundary is slightly less precise than the CNN's learned representations. The BiLSTM is
the clear outlier at 0.789, confirming that its weaker performance reflects genuinely
poor class separation in probability space rather than a miscalibrated decision threshold
--- the recurrent model simply fails to learn discriminative representations from MFCC
sequences for this task.

\begin{figure}[h!]
    \centering
    \includegraphics[width=0.85\textwidth]{figures/roc_curves.png}
    \caption{Macro-averaged one-vs-rest ROC curves for all three models on the 2{,}645
             test samples. Each curve averages 12 per-class ROC curves. The CNN achieves
             the highest macro AUC (0.893), with all three models substantially above the
             random baseline (dashed line, AUC = 0.500).}
    \label{fig:roc}
\end{figure}

% ─────────────────────────────────────────────────────────────────────────────
\section{Discussion and Limitations}
\label{sec:discussion}
% ─────────────────────────────────────────────────────────────────────────────

\subsection{CNN vs.\ BiLSTM: Spectral Texture vs.\ Temporal Dynamics}

The CNN's superior performance (50.59\% vs.\ 40.42\%) suggests that spectral texture
captured in a Mel spectrogram image is more discriminative for genre classification than
temporal dynamics captured in MFCC sequences. Music genre is primarily characterised by
timbre, instrumentation, and energy distribution --- properties well-encoded in a
time-frequency image --- rather than by how those properties evolve over exactly 130
short-term frames. The CNN can detect local spectral patterns (e.g.\ the dense
broadband energy of Electronic or the sparse harmonic structure of Classical) directly
from the image, whereas the BiLSTM must infer genre from the statistical evolution of
cepstral coefficients, which is a more indirect and lossy representation.

\subsection{Competitive SVM Baseline}

The SVM baseline (48.05\%) performs surprisingly close to the CNN, outperforming the
BiLSTM by nearly 8 percentage points. This aligns with prior work \citep{tzanetakis2002musical},
which demonstrated that MFCC summary statistics carry substantial genre-discriminative
information. The SVM's competitiveness highlights that the BiLSTM's deficit is not
simply due to using MFCC features --- it is the sequential formulation that underperforms,
not the features themselves. It should be noted that the 2.5 percentage point accuracy
gap between the CNN and SVM has not been subjected to a statistical significance test;
a McNemar's test on the paired test-set predictions would be required to confirm whether
this difference is statistically meaningful rather than within the range of chance variation.

\subsection{Hard and Easy Genres}

\textbf{Easy genres:} Hip-Hop (CNN F1: 0.762) and Classical (CNN F1: 0.714) are the
most reliably classified across all models. Hip-Hop has strong low-frequency
periodicity (bass and beat patterns) that is acoustically distinct; Classical exhibits
sparse, structured harmonic content with wide dynamic range --- both are visually
separable in the Mel spectrogram space.

\textbf{Hard genres:} Pop (CNN F1: 0.168), Jazz (CNN F1: 0.265), and Instrumental
(CNN F1: 0.354) are consistently the weakest performers. Pop is a stylistically
heterogeneous genre that borrows elements from Rock, Electronic, Hip-Hop, and Folk,
making it acoustically ill-defined. Jazz is a small class (only 46 test samples) with
high acoustic overlap with Instrumental and Classical. Instrumental, by definition,
spans many genres without lyrics as a discriminating signal, making it difficult to
distinguish from Rock, Folk, or Electronic depending on the clip.

\subsection{Overfitting Analysis}

Both deep models show a train/validation accuracy gap, consistent with the
moderate overfitting expected on a dataset of this size. Dropout (0.5 and 0.3),
gradient clipping, and early stopping partially mitigate this. The BiLSTM's
training curves plateau earlier than the CNN's, suggesting the recurrent model
reaches its representational capacity faster on this task.

\subsection{Limitations and Future Work}

\begin{itemize}[noitemsep, topsep=4pt]
    \item \textbf{Genre labels are subjective.} Many real-world tracks span multiple
          genres; single-label classification imposes a hard boundary that may not
          reflect the underlying acoustic reality.
    \item \textbf{Class imbalance persists.} Despite filtering and class-weighted loss,
          Rock and Electronic together account for $\approx$35\% of clips, which likely
          biases all models toward these majority classes.
    \item \textbf{Fixed clip duration.} Truncating or padding to 30 seconds discards
          long-range musical structure and may cut clips mid-phrase.
    \item \textbf{No data augmentation.} Time-stretching, pitch-shifting, or mixup
          augmentation could improve generalisation, particularly for underrepresented
          genres such as Country and Jazz.
\end{itemize}

The most promising future extension would be a \textbf{Convolutional Recurrent Neural
Network (CRNN)} \citep{choi2017crnn} combining both branches: CNN convolutional blocks
extract local spectral features from the Mel spectrogram, which are then fed as a
sequence to an LSTM for temporal modelling. This architecture directly addresses the
complementary weaknesses of our two models. Transfer learning from a large pretrained
audio model (e.g.\ VGGish trained on AudioSet) is another viable direction, particularly
for the small classes where our models struggle most.

% ─────────────────────────────────────────────────────────────────────────────
\section{Conclusion}
\label{sec:conclusion}
% ─────────────────────────────────────────────────────────────────────────────

This report presented a complete deep learning pipeline for automatic music genre
classification using the \texttt{lewtun/music\_genres} dataset. After exploratory data
analysis and justified filtering from 20 to 12 genres, two complementary architectures
were trained and evaluated: a 2-D CNN operating on log-scaled Mel spectrograms, and a
bidirectional LSTM operating on MFCC frame sequences, alongside a classical SVM baseline.
The CNN achieved the highest test accuracy (50.59\%, weighted F1: 0.507), outperforming
the SVM (48.05\%) and the BiLSTM (40.42\%), indicating that spectral texture captured in
a time-frequency image is more discriminative for genre than temporal dynamics captured
in MFCC sequences alone. Hip-Hop and Classical were consistently the easiest genres to
classify, while Pop, Jazz, and Instrumental proved difficult across all models due to
acoustic heterogeneity and class overlap. Future work should explore a CRNN architecture
that unifies both representation types, as well as data augmentation strategies and
transfer learning from large pretrained audio models to improve performance on
underrepresented genres.

% ─────────────────────────────────────────────────────────────────────────────
\newpage
\bibliographystyle{plainnat}

\begin{thebibliography}{99}

\bibitem[Choi et~al.(2017a)]{choi2017convolutional}
Choi, K., Fazekas, G., Sandler, M., and Cho, K. (2017a).
\newblock Convolutional recurrent neural networks for music classification.
\newblock In \textit{IEEE International Conference on Acoustics, Speech and Signal
Processing (ICASSP)}.

\bibitem[Choi et~al.(2017b)]{choi2017crnn}
Choi, K., Fazekas, G., and Sandler, M. (2017b).
\newblock Music tagging with deep recurrent neural networks.
\newblock In \textit{Proceedings of ISMIR 2017}.

\bibitem[Hershey et~al.(2017)]{hershey2017cnn}
Hershey, S., Chaudhuri, S., Ellis, D.~P., et~al. (2017).
\newblock {CNN} architectures for large-scale audio classification.
\newblock In \textit{IEEE International Conference on Acoustics, Speech and Signal
Processing (ICASSP)}.

\bibitem[LeCun et~al.(2015)]{lecun2015deep}
LeCun, Y., Bengio, Y., and Hinton, G. (2015).
\newblock Deep learning.
\newblock \textit{Nature}, 521(7553):436--444.

\bibitem[Kingma and Ba(2014)]{kingma2014adam}
Kingma, D.~P. and Ba, J. (2014).
\newblock Adam: A method for stochastic optimization.
\newblock \textit{arXiv preprint arXiv:1412.6980}.

\bibitem[Srivastava et~al.(2014)]{srivastava2014dropout}
Srivastava, N., Hinton, G., Krizhevsky, A., Sutskever, I., and Salakhutdinov, R. (2014).
\newblock Dropout: A simple way to prevent neural networks from overfitting.
\newblock \textit{Journal of Machine Learning Research}, 15(1):1929--1958.

\bibitem[Tzanetakis and Cook(2002)]{tzanetakis2002musical}
Tzanetakis, G. and Cook, P. (2002).
\newblock Musical genre classification of audio signals.
\newblock \textit{IEEE Transactions on Speech and Audio Processing}, 10(5):293--302.

\end{thebibliography}

% ─────────────────────────────────────────────────────────────────────────────
\appendix
\section{Code Repository and Implementation Details}
\label{sec:appendix}
% ─────────────────────────────────────────────────────────────────────────────

The full source code, preprocessing notebooks, and model training notebooks are
available at:

\begin{center}
\url{https://github.com/marticasasn/ELEC3612-Assignment2}
\end{center}

The repository includes:
\begin{itemize}[noitemsep, topsep=4pt]
    \item \texttt{notebooks/01\_eda\_kaggle.ipynb} --- exploratory data analysis and
          genre filtering
    \item \texttt{notebooks/02\_preprocessing\_kaggle.ipynb} --- feature extraction
          and tensor saving
    \item \texttt{notebooks/03\_cnn\_kaggle.ipynb} --- CNN training and evaluation
    \item \texttt{notebooks/04\_lstm\_kaggle.ipynb} --- BiLSTM training and evaluation
    \item \texttt{notebooks/05\_evaluation\_kaggle.ipynb} --- comparative evaluation
          and SVM baseline
    \item \texttt{src/} --- reusable PyTorch modules: dataset loaders, model
          architectures, training loop, and evaluation utilities
\end{itemize}

\end{document}